% Nouveau EGA V (cote 25): archive cote 25, site page 97 (PDF page 98).
% typescript. Draft compared: none.
% AI-transcribed by Claude Opus 5.5 (claude-opus-5-5, adaptive thinking, effort high), Message Batches API, batch ega25b (msgbatch_01A7n5YuXubUPNK1zsepi8RG), 2026-11-13; not yet checked by eye.
\scanpage{25/97}
% NOTE: the whole typed text of this page is extremely faint (pale carbon or show-through); most of it can only be guessed, and many readings below are uncertain
% NOTE: two or three faint typed lines at the head of the page, above the heading, cannot be read
\illegible{8}

\illegible{8}

\uline{\unsure{§ 2)} \unsure{Compléments sur les schémas projectifs} \illegible{3}} \unsure{et}

\uline{\unsure{Pic(X). Les} \illegible{3}}

\uline{\unsure{1. Détermination des faisceaux inversibles sur les fibrés projectifs.}}

\uline{\unsure{a) Fibrés en espaces projectifs associés à un fibré projectif.}}

\uline{\unsure{Proposition 1.1.}} \unsure{--- Soient S un préschéma, E un Module loc. libre sur S partout}
\unsure{(rang $\geq 2$ en tout point).}
\xout{\illegible{3}}\unsure{, $X = P(\underline{E})$ le fibré projectif qu'il définit. Alors}
\unsure{pour tout Module inversible $\underline{L}$ sur X, on peut trouver une famille}
\unsure{finie $(S_i)_{i \in I}$ d'ouverts disjoints de S recouvrant S, indexés par I, et}
\unsure{pour tout $i \in I$ un Module inversible $\underline{M}_i$ sur $S_i$, et entiers $n_i$ tels que}
\unsure{la rest.} \xout{\illegible{4}} \unsure{la rest.}
\unsure{de $\underline{L}$ à $f^{-1}(S_i)$ soit isomorphe à celle de $f^*(\underline{M}_i)(n_i)$.}
\unsure{De plus, la famille des $S_i$} \xout{\illegible{2}} \unsure{est uniquement déter-}%
\unsure{minée par ces conditions, ainsi que $\underline{M}_i$ à isomorphisme unique près.}
\unsure{Corollaire 1.2.} \unsure{--- Si on ne fait pas d'hypothèse sur le rang de E, S se}
\unsure{décompose canoniquement en somme de trois ouverts $S'$, $S''$, $S'''$ tels que sur}
\unsure{$S'$, $S''$, $S'''$, le rang de $\underline{E}$ soit 0, 1, $\geq 2$. Alors la détermination des}
\unsure{Modules inversibles sur X se ramène à celle des $f^{-1}(S''')$ par 1.1.1.}
\unsure{Le cas 1.2 est justiciable de 1.1. D'autre part $f^{-1}(S'')$ est isomorphe}
\unsure{à $S''$. Donc son groupe de Picard n'est autre que $\mathrm{Pic}(S'')$, enfin}
\unsure{$f^{-1}(S')$ est vide, donc son groupe de Picard est nul.}
\uline{\unsure{Corollaire 1.3.}} \unsure{Sous les conditions de 1.1, supposons S connexe. Alors}
\unsure{tout Module inversible $\underline{L}$ sur X est isomorphe à un Module de la forme}
\unsure{$f^*(\underline{M})(n)$,} \illegible{2} \unsure{$n \in \mathbb{Z}$,} \illegible{3} \unsure{de plus il est}
\illegible{6} \unsure{de la forme}

\unsure{$*$}
\unsure{Un autre} \illegible{8}
\illegible{6}

\mspage{\hand{97}}
